% TOP_PROBLEM: {"id": 3146, "title": "Half-integral flow polynomial values", "category_id": 3, "category": {"id": 3, "name": "graph_theory", "display_name": "Graph Theory", "description": "Problems involving graphs, networks, and their properties.", "slug": "graph-theory", "order_index": 3, "created_at": "2026-07-31T15:26:25.671Z"}, "status": "open", "problem_number": "OPG-348", "set_id": 12, "statusQualification": "Flow-polynomial counts are specified as finite-group (modular) flows, resolving the source shorthand “number of nowhere-zero k-flows.”"}
% FIRST_POSTED: 2026-10-01
% ENTRY_KIND: statement_only
% UnsolvedMath ID 3146; source catalogue code OPG-348
% !TeX program = lualatex
% Definition and sources only; this file is not an LLM solution attempt.
\documentclass[11pt,a4paper]{article}
\usepackage[margin=25mm]{geometry}
\usepackage{fontspec}
\setmainfont{lmroman10-regular.otf}[BoldFont=lmroman10-bold.otf,
 ItalicFont=lmroman10-italic.otf,BoldItalicFont=lmroman10-bolditalic.otf]
\usepackage{amsmath,amssymb,mathtools}
\usepackage{unicode-math}
\setmathfont{latinmodern-math.otf}
\AtBeginDocument{\let\setminus\smallsetminus}
\usepackage{xurl,hyperref}
\hypersetup{colorlinks=true,urlcolor=blue}
\setcounter{secnumdepth}{0}
\setlength{\parindent}{0pt}
\setlength{\parskip}{5pt}
\setlength{\emergencystretch}{3em}
\begin{document}
\section*{Half-integral flow polynomial values}\label{3146}
{\small UnsolvedMath ID 3146 · OPG-348 · Graph Theory · OpenGarden}
\subsection{Definitions and mathematical statement}
Let $G=(V,E)$ be a finite connected loopless undirected multigraph;
parallel edges are regarded as distinct. Assume $G$ has no bridge,
where a bridge is an edge whose deletion disconnects the graph.
Choose an orientation of its edges. For a positive integer $q$, a
nowhere-zero $\mathbb Z/q\mathbb Z$-flow assigns a nonzero residue
$\phi(e)$ to every edge so that the sum entering each vertex equals
the sum leaving it, with equality in $\mathbb Z/q\mathbb Z$.

The flow polynomial $\Phi(G,x)$ is characterized by counting these
flows at every positive integer $q$. A definition valid directly
for all real $x$ is
\[
 \Phi(G,x)=\sum_{A\subseteq E}(-1)^{|E|-|A|}
 x^{|A|-|V|+c(A)},
\]
where $c(A)$ is the number of connected components of the spanning
graph $(V,A)$, including isolated vertices. Each exponent is
nonnegative. Reversing an edge and negating its value gives a
bijection of flow sets, so the count is independent of orientation.

The conjecture asks whether every such $G$ satisfies
\[
 \Phi(G,11/2)>0.
\]
The argument $11/2$ is a polynomial evaluation, not the order of
a finite group. In particular, positivity at neighbouring integer
arguments does not by itself give the requested inequality.
The counting convention uses modular flows; counting bounded
integer-valued flows instead would define a different counting
function and would not specify this polynomial.
\subsection{Short English statement}
Is the flow polynomial of every bridgeless connected graph strictly positive at five and a half?
\subsection{Sources}
\begin{itemize}
\item[S1] ULAM.AI and the UnsolvedMath Contributors, supplied problems.json, record 3146 (OPG-348), OpenGarden collection. Catalogue identity and source transcription; CC BY 4.0.
\url{https://huggingface.co/datasets/ulamai/UnsolvedMath}.
\item[S2] Open Problem Garden, Half-integral flow polynomial values. Original problem and discussion; the mathematical formulation below is expanded from the supplied source record.
\url{http://www.openproblemgarden.org/op/half_integral_flow_polynomial_values}.
\end{itemize}
\end{document}
